\documentclass[12pt]{beamer}

\usetheme{Boadilla}
\usecolortheme{default}

\usepackage[utf8]{inputenc}
\usepackage[T1]{fontenc}

\usepackage{amsmath, amssymb}
\usepackage{booktabs}

% Informações
\title{Por que precisamos de Ecologia?}
\subtitle{E por que a matemática é sua linguagem natural}
\author{Modelagem Matemática}
\institute{Matemática Aplicada}
\date{Aula 01}

\begin{document}

% ── CAPA ──────────────────────────────────────────
\begin{frame}
  \titlepage
\end{frame}



% ══════════════════════════════════════════════════
\section{Um colapso que ninguém previu}
% ══════════════════════════════════════════════════

\begin{frame}{Um colapso que ninguém previu}

  Em 1989, a população de bacalhau do Atlântico Norte entrou em colapso
  após \textbf{500 anos} de pesca contínua.
  Em 1992, o governo canadense decretou \textbf{moratória total}.

  \bigskip
  Trinta anos depois --- a população \textbf{não se recuperou}.

  \bigskip
  \begin{block}{Paradoxo}
    Pescadores, gestores e biólogos observavam o mesmo sistema
    há gerações. Tinham intuição, experiência e dados de captura.
    E ainda assim --- não previram, não evitaram, não reverteram.
  \end{block}

  \bigskip
  \centering\large\itshape Por quê?

\end{frame}

% ── Slide 1/2: dados do Pantanal — análogo brasileiro ─────────────────────────
\begin{frame}{O mesmo problema acontece no Brasil}

  \mbox{\footnotesize CPUE — Pantanal MS, 1994--2017 \textbullet{} SCPESCA/MS, Embrapa Pantanal}

  \vspace{3pt}
  \begin{columns}[T]
    \begin{column}{0.60\textwidth}
      \includegraphics[width=\linewidth]{cpue_pantanal.pdf}
    \end{column}
    \begin{column}{0.37\textwidth}
      \textbf{O que os dados mostram:}
      {\footnotesize
      \begin{itemize}\setlength\itemsep{3pt}\setlength\parskip{0pt}
        \item \textbf{Pacu}: CPUE caiu de 14 para 2{,}5\,kg/pescador em 6 anos com esforço crescente
        \item Após restrição (2000): patamar \textbf{permanentemente inferior}
        \item \textbf{Pintado} e \textbf{dourado}: declínio contínuo
      \end{itemize}
      \vspace{4pt}
      CPUE é o proxy de abundância do estoque.}
    \end{column}
  \end{columns}

\end{frame}

% ── Slide 2/2: confronto bacalhau × Pantanal ──────────────────────────────────
\begin{frame}{Bacalhau e Pantanal: o mesmo argumento, dois contextos}

  \begin{block}{Padrão comum}
    {\footnotesize\textbf{(1)} Captura cai com esforço crescente\enspace
    \textbf{(2)} Dados contínuos não previram o declínio\enspace
    \textbf{(3)} Linguagem verbal não distingue recuperação de colapso}
  \end{block}

  \vspace{3pt}
  \begin{columns}[T]
    \begin{column}{0.47\textwidth}
      \begin{exampleblock}{Bacalhau — Atlântico N.}
        {\footnotesize
        \begin{itemize}\setlength\itemsep{0pt}\setlength\parskip{0pt}
          \item 500 anos de pesca contínua
          \item Colapso 1989, moratória 1992
          \item Não se recuperou em 30 anos
        \end{itemize}}
      \end{exampleblock}
    \end{column}
    \begin{column}{0.47\textwidth}
      \begin{exampleblock}{Pescados — Pantanal MS}
        {\footnotesize
        \begin{itemize}\setlength\itemsep{0pt}\setlength\parskip{0pt}
          \item 23 anos de dados (SCPESCA)
          \item Declínio detectado; restrições aplicadas
          \item Recuperação parcial e lenta
        \end{itemize}}
      \end{exampleblock}
    \end{column}
  \end{columns}

  \vspace{3pt}
  \begin{alertblock}{A pergunta de Krebs, reformulada}
    {\footnotesize Quais \textbf{interações} entre esforço de pesca,
    reprodução e abundância determinam \emph{onde} o sistema vai parar?}
  \end{alertblock}

\end{frame}

% ══════════════════════════════════════════════════
\section{O limite da intuição verbal}
% ══════════════════════════════════════════════════

\begin{frame}{O limite da intuição verbal}

  Considere estas afirmações --- todas razoáveis, todas verbais:

  \bigskip
  \begin{enumerate}
    \item ``Se pescamos menos, o estoque se recupera.''
    \item ``A população é resiliente --- sempre se recuperou antes.''
    \item ``O ambiente mudou, por isso não volta.''
  \end{enumerate}

  \bigskip
  \begin{alertblock}{Problema}
    Todas as três podem ser \textbf{verdadeiras}. Ou \textbf{falsas}.
    Ou verdadeiras sob certas condições e falsas sob outras.
    A linguagem verbal não tem como distinguir entre elas:
    ela não diz \emph{quando}, \emph{por que} nem \emph{sob quais condições}.
  \end{alertblock}

\end{frame}

% ══════════════════════════════════════════════════
\section{Por que precisamos de ecologia?}
% ══════════════════════════════════════════════════

\begin{frame}{Por que precisamos de ecologia?}

  Observação bruta descreve \textbf{padrões}.
  Ecologia explica \textbf{mecanismos}.

  \bigskip
  \begin{block}{Observação}
    \begin{itemize}
      \item Captura caiu 90\% em 10 anos
      \item Estoque não volta após moratória
      \item Temperatura do oceano subiu
    \end{itemize}
  \end{block}

  \begin{exampleblock}{Ecologia pergunta}
    \begin{itemize}
      \item Quais interações causaram o colapso?
      \item Existe um ponto sem retorno?
      \item Como o ambiente altera a dinâmica?
    \end{itemize}
  \end{exampleblock}

\end{frame}

% ══════════════════════════════════════════════════
\section{Conclusão 1 — A pergunta de Krebs}
% ══════════════════════════════════════════════════

\begin{frame}{Conclusão 1 — A pergunta que organiza tudo}

  Chegamos a uma definição que não é ponto de partida --- é \textbf{resposta}:

  \bigskip
  \begin{block}{Definição de Krebs (1972)}
    \itshape
    ``A ecologia é o estudo científico das \textbf{interações}
    que \textbf{determinam} a \textbf{distribuição e a abundância}
    dos organismos.''
  \end{block}

  \vspace{0.3em}
  {\small Adotada por Begon, Townsend \& Harper (2006).}

  \bigskip
  Reformulada como pergunta:

  \bigskip
  \begin{exampleblock}{A pergunta central}
    \centering
    Quais interações determinam \emph{onde} os organismos
    ocorrem e em \emph{que quantidade}?
  \end{exampleblock}

\end{frame}

% ── ANATOMIA DA DEFINIÇÃO ─────────────────────────
\begin{frame}{Três palavras que carregam todo o peso}

  \begin{itemize}
    \item \textbf{Interações} --- relação de dependência entre variáveis
      que mudam no tempo
    \bigskip
    \item \textbf{Determinam} --- causalidade: não basta descrever,
      é preciso explicar
    \bigskip
    \item \textbf{Distribuição e abundância} --- grandezas mensuráveis:
      o objeto já é quantitativo
  \end{itemize}

  \bigskip
  \begin{exampleblock}{Implicação}
    A ecologia não matematiza a biologia por opção estética.
    Seu objeto central --- abundância variável no tempo ---
    \textbf{já é matemático}.
  \end{exampleblock}

\end{frame}

% ══════════════════════════════════════════════════
\section{Por que precisamos de modelos?}
% ══════════════════════════════════════════════════

\begin{frame}{Por que precisamos de modelos?}

  A pergunta de Krebs exige mais do que observação e descrição qualitativa.
  Precisamos de um objeto que:

  \bigskip
  \begin{itemize}
    \item force \textbf{precisão nas hipóteses}
    \item gere \textbf{predições testáveis}
  \end{itemize}

  \bigskip
  Esse objeto é o \textbf{modelo matemático}.

  \bigskip
  \begin{alertblock}{Questão}
    Mas o que construir um modelo, de fato, nos força a fazer?
  \end{alertblock}

\end{frame}

% ── O QUE O MODELO EXIGE ──────────────────────────
\begin{frame}{O que construir um modelo nos força a responder}

  \begin{itemize}
    \item \textbf{A que taxa os indivíduos nascem e morrem?}\\
      {\small $\rightarrow$ A dinâmica é o mecanismo, não o padrão}
    \bigskip
    \item \textbf{Essa taxa depende do tamanho atual da população?}\\
      {\small $\rightarrow$ Dependência de densidade --- existe um limite}
    \bigskip
    \item \textbf{Como a pesca altera essa taxa?}\\
      {\small $\rightarrow$ A intervenção entra como parâmetro}
    \bigskip
    \item \textbf{O sistema tem mais de um equilíbrio?}\\
      {\small $\rightarrow$ Pode haver um ponto sem retorno}
  \end{itemize}

\end{frame}

% ══════════════════════════════════════════════════
\section{Conclusão 2 — A linguagem necessária}
% ══════════════════════════════════════════════════

\begin{frame}{Conclusão 2 — A linguagem necessária}

  Para responder à pergunta de Krebs, precisamos representar como
  \textbf{quantidades mudam em função de interações}.

  \bigskip
  O modelo mais simples possível de uma população:

  \[
    \frac{dN}{dt} = f(N)
  \]

  \centering{\small taxa de mudança da abundância $=$ função das interações}

  \bigskip
  \begin{exampleblock}{Conclusão}
    A linguagem que representa mudança por interação é o
    \textbf{cálculo diferencial}. A ecologia chega à matemática
    não por abstração --- mas por \textbf{necessidade}.
  \end{exampleblock}

\end{frame}

% ══════════════════════════════════════════════════
\section{Síntese}
% ══════════════════════════════════════════════════

\begin{frame}{Síntese — O fio condutor}

  \begin{enumerate}
    \item \textbf{Problema concreto} \hfill colapso do bacalhau
    \bigskip
    \item \textbf{Intuição verbal falha} \hfill precisamos de mecanismo
    \bigskip
    \item \textbf{Conclusão 1} \hfill a pergunta de Krebs
    \bigskip
    \item \textbf{Descrição qualitativa falha} \hfill precisamos de modelo
    \bigskip
    \item \textbf{Conclusão 2} \hfill cálculo diferencial como linguagem
  \end{enumerate}

\end{frame}

% ── PRÓXIMA AULA ──────────────────────────────────
\begin{frame}{Próxima aula — Da pergunta ao primeiro modelo}

  Com a linguagem estabelecida, construímos o modelo mais simples
  que responde à pergunta de Krebs para uma única espécie:

  \[
    \frac{dN}{dt} = rN\!\left(1 - \frac{N}{K}\right)
  \]

  Onde cada símbolo é uma \textbf{hipótese explícita} sobre o mecanismo:

  \bigskip
  \begin{itemize}
    \item $N$ --- abundância (o que Krebs quer explicar)
    \item $r$ --- taxa intrínseca de crescimento (interação indivíduo-recurso)
    \item $K$ --- capacidade de suporte (limite imposto pelo ambiente)
  \end{itemize}

\end{frame}

\end{document}
